\section{Experimental setup}
\label{sec:setup}

\paragraph{Agents.} We evaluate commercial agents as they are shipped, that is,
a model together with the harness that drives it. We call each such pairing an
\emph{arm}. We do not have API-level access that would let us hold the
scaffold fixed and swap only the model, so every claim in this paper is about
arms, not about raw models. One arm runs Claude Opus~4.8 through the Claude
Code command-line interface. All other arms run through \texttt{cursor-agent},
the headless command-line client of the Cursor editor, which exposes many
vendors' models behind a single scaffold. All Cursor arms are invoked through the same client
with the same flags and differ only in the model flag, so differences among
them are closer to differences between models than they would be across
vendors' own products. We did not inspect the client's internal prompts and
cannot rule out that it adapts them per model.

We ran two generations. \emph{Generation~1} (September 2026) covers eight arms
on the seven core domains. \emph{Generation~2} (late September and October
2026) covers the newest model available from each vendor through Cursor, on
the full suite. Table~\ref{tab:armlist} lists the arms with the exact model
identifiers passed to the client. Three identifiers (Composer~2.5, GLM~5.2,
and GPT-5.6~Sol) appear in both generations unchanged, and Codex~5.3 appears
with a different reasoning setting, so the \NArms{} arm-generation pairs we
analyze correspond to \NConfigsDistinct{} distinct model-and-harness
configurations.

\begin{table}[t]
\centering
\footnotesize
\caption{Agent arms. All Cursor arms share the same client and scaffold.
Model identifiers are the strings passed to the client's model flag.}
\label{tab:armlist}
\begin{tabular}{@{}llll@{}}
\toprule
Generation & Vendor & Harness & Model identifier \\
\midrule
1 & Anthropic & Claude Code & \texttt{claude-opus-4-8} \\
1 & Cursor    & cursor-agent & \texttt{composer-2.5} \\
1 & Moonshot  & cursor-agent & \texttt{kimi-k2.7-code} \\
1 & Zhipu     & cursor-agent & \texttt{glm-5.2-high} \\
1 & OpenAI    & cursor-agent & \texttt{gpt-5.6-sol-high} \\
1 & OpenAI    & cursor-agent & \texttt{gpt-5.3-codex} \\
1 & xAI       & cursor-agent & \texttt{cursor-grok-4.6-high} \\
1 & Google    & cursor-agent & \texttt{gemini-3.6-flash-high} \\
\midrule
2 & Anthropic & cursor-agent & \texttt{claude-opus-5-5-high} \\
2 & Cursor    & cursor-agent & \texttt{composer-2.5} \\
2 & Moonshot  & cursor-agent & \texttt{kimi-k3-high} \\
2 & Zhipu     & cursor-agent & \texttt{glm-5.2-high} \\
2 & OpenAI    & cursor-agent & \texttt{gpt-5.6-sol-high} \\
2 & OpenAI    & cursor-agent & \texttt{gpt-5.3-codex-high} \\
2 & xAI       & cursor-agent & \texttt{grok-4.7-high} \\
2 & Google    & cursor-agent & \texttt{gemini-3.8-flash-high} \\
\bottomrule
\end{tabular}
\end{table}

\paragraph{Protocol.} For each run the harness creates a fresh directory,
copies in the task repository and the visible tests only, starts the agent
in that directory, and invokes it non-interactively with the contents of \texttt{prompt.md}. Claude Code
runs with permission prompts disabled and a USD~2 budget cap; Cursor runs
with the flags \texttt{--print} and \texttt{--force}. The harness sets a 600-second time limit on the client process
(\S\ref{sec:limits} discusses runs that exceeded it). No
human intervenes between the prompt and the agent's exit. The full event
stream is saved as the run's trajectory. When the agent exits, the harness
restores the hidden tests into a separate copy, runs both test sets with
\texttt{pytest}, computes the line and symbol diffs against the pristine
repository, and writes one row per run to a SQLite database. We measure
exploration as \emph{unique tool calls}: the number of distinct tool-call
identifiers in a run's trajectory. The two harnesses record tool calls in
different formats, so these counts are comparable within a harness, and we
use them only for within-harness contrasts. Claude Code reports cost and the
Cursor stream does not, so we do not report cost.

\paragraph{Invalid runs.} A run is valid if the harness recorded no error and
the trajectory contains at least one tool call. The harness records an error
when the client returns a vendor usage-cap notice, when the session ends with
no tool calls and a transcript too short to contain an attempt, or when the
run hits the time limit. Invalid runs are excluded from every rate.
Without this rule such a session would be scored as a clean zero-drift run,
since an untouched repository has no unrequested changes. Invalid runs stay in
the database with the reason recorded, and Table~\ref{tab:coverage} reports
how many each arm lost.

\paragraph{Statistics.} Most cells are a single run ($k=1$), because of the
usage limits of the subscriptions we had. We report proportions with Wilson
95\% intervals, which treat runs as independent and should be read as
descriptive, since runs on the same domain are related. Differences and means
use a task-level bootstrap (10{,}000 resamples, fixed seed). The ladder effect is tested with Cochran's $Q$ on complete three-rung
blocks (one arm on one domain), with the $p$-value taken from 10{,}000
permutations of rung labels within domain, applied identically to every arm
of that domain, so that agents solving the same domain are not treated as
independent. Since the same repository appears under all three rungs, we also
report paired comparisons within arm and domain.
